\begin{itemize}
\item \textbf{Vehicle.} The booster is a rigid circular cylinder, 65 m long and 7 m in diameter. Its 130 t dry mass is distributed uniformly for the pitch inertia, and the remaining propellant is treated as a point mass at the center of mass (Section~\ref{sec:motion}). Pitch rotation is modeled because the offset between the center of mass and the center of pressure creates aerodynamic moments, so an upright touchdown has to follow from the dynamics. Structural and thermal capacity are outside the study; the loads are reported as requirements.
\item \textbf{Earth and atmosphere.} The reference calculation uses a nonrotating spherical Earth with inverse-square gravity and the 1976 U.S. Standard Atmosphere, continued isothermally above 84.852 km geopotential altitude where needed. A flat-Earth version of the same equations provides starting guesses and the earlier planar comparisons. Motion stays in the vertical plane of the separation velocity, and the nominal case has no wind.
\item \textbf{Separation state.} The booster separates at 70 km altitude and 6000 km/h. The flight-path angle, which the problem does not give, is $15^\circ$ above horizontal in the reference case and $0^\circ$ or $30^\circ$ in the alternatives. The body axis is initially aligned with the velocity and the pitch rate is zero.
\item \textbf{Propulsion.} One engine provides up to 2.846 MN, throttleable between 40\% and 100\%, with a $\pm10^\circ$ gimbal, an effective specific impulse of 330 s, and an oxygen-to-methane mass ratio of 3.5, so methane is $1/4.5$ of the propellant mass. At most three ignitions are allowed. Throttle and gimbal respond instantaneously, with no ignition delay or minimum burn time. There are no fins or reaction-control thrusters, so the attitude evolves passively while the engine is off.
\item \textbf{Aerodynamics.} Axial and crossflow forces use constant coefficients $C_a=0.8$ and $C_n=1.2$, with the center of pressure $0.05L$ ahead of the center of mass.
\item \textbf{Ship and touchdown.} The ship is stationary, with a level deck at sea level, and may be placed anywhere along the ground track; each scenario gets its own position. The booster has no landing gear. A landing is accepted when, at first contact, the booster is within 5 m of the ship's position, moves less than 1 m/s horizontally and 2 m/s vertically, tilts less than $2^\circ$ from upright, rotates slower than $1^\circ$/s, and retains at least 5 t of usable propellant. Unusable residual propellant is taken as zero.
\item \textbf{Parameter studies.} We vary the separation angle ($0^\circ$, $15^\circ$, $30^\circ$), the effective specific impulse (300, 330, 360 s), and the crossflow coefficient (0.8, 1.0, 1.2, 1.6), one at a time. The axial coefficient, pressure-center offset, thrust, and reserve are not varied.
\end{itemize}
